\subsection*{4.3 Distance-efficient rate amplification}
Let $P$ be a set, and $\mathcal R$ a partition of $P^2$. We denote the part containing $p$ by $[p]_{\mathcal R}$ or $[p]$ when $\mathcal R$ is clear from context. We call $(p)=[p]\setminus\{p\}$ the open class of $p$. For a set $A\subseteq P^2$ we let $(A)=\bigcup_{p\in A}(p)$. Given $p\in P$ we say that $\{p\}\times P\subseteq P^2$ is vertical line and $P\times\{p\}$ is a horizontal line. Horizontal and vertical lines are referred to as axis-parallel lines, and we denote the set of such lines by
\[\mathcal X=\bigcup_{p\in P}\{\{p\}\times P,P\times\{p\}\}.\]
For a point $p=(p_1,p_2)\in P^2$ we denote $S_p=(\{p_1\}\times P)\cup(P\times\{p_2\})\setminus\{p\}$. That is, $S_p$ is the set of points in $P^2$ of distance exactly $1$ from $p$. Key to our distance-efficient rate amplification procedure is a partition of the ``square'' $P^2$ with certain properties.
\setcounter{nttheorem}{34}
\begin{ntdefinition}[Author Definition35: Axis-evasive partitions]
Let $P$ be a set. A partition $\mathcal R$ of $P^2$ is said to be $(c,s)$-axis evasive if
\begin{enumerate}
\item For every $p\in P^2$, $|(p)|\leq c$.
\item For every $\ell,\ell'\in\mathcal X$ (possibly equal), $|\ell'\cap(\ell)|\leq s$.
\item For every $p\in P^2$ and $\ell\in\mathcal X$, $|[p]\cap\ell|\leq1$.
\end{enumerate}
\end{ntdefinition}
In Section5 we study such partitions. We prove their existence with certain parameters and give explicit constructions. In this section, however, we work with abstract axis-evasive partitions and analyze our rate amplification procedure with respect to the parameters $c,s$ of the axis-evasive partition as well as the number of parts which we typically denote by $t$.
\begin{ntclaim}[Author Claim36]
Let $p,p'\in P^2$ (possibly equal). Then, $|\{r\in S_p\mid(r)\cap S_{p'}\ne\varnothing\}|\leq4s$.
\end{ntclaim}
\begin{proof}
Note that each of $S_p,S_{p'}$ is a subset of the union of two axis-parallel lines. Thus, to prove the claim, it suffices to show that for every $\ell,\ell'\in\mathcal X$, not necessarily distinct,
\[|\{r\in\ell\mid(r)\cap\ell'\ne\varnothing\}|\leq s.\]
Let $r_1,\ldots,r_t\in\ell$ be such that $(r_i)\cap\ell'\ne\varnothing$. Note that for every distinct $i,j\in[t]$ it holds that $((r_i)\cap\ell')\cap((r_j)\cap\ell')=\varnothing$. Indeed, since $\mathcal R$ is a partition, if $((r_i)\cap\ell')\cap((r_j)\cap\ell')\ne\varnothing$ then $r_i\in[r_j]$, but this implies that $|\ell\cap[r_j]|\geq2$ in contradiction axis evasiveness. Thus,
\[R=\bigcup_{i=1}^t((r_i)\cap\ell')\]
is a disjoint union of size $t$. However, $R\subseteq(\ell)\cap\ell'$, and so $t\leq|R|\leq|(\ell)\cap\ell'|\leq s$.
\end{proof}
\begin{ntdefinition}[Author Definition37]
Let $P$ be a set, $\mathbb F$ a field. Let $\mathcal R$ be a $(c,s)$-axis evasive partition of $P^2$. For every $p\in P^2$ define the function $g_{[p]}:P^2\to\mathbb F$ as follows:
\[g_{[p]}(r)=\begin{cases}1,&r\in[p];\\0,&\text{otherwise}.\end{cases}\]
We define $L_{\mathcal R}=\{g_{[p]}\mid p\in P^2\}$.
\end{ntdefinition}
\begin{ntdefinition}[Author Definition38]
Let $P$ be a set, $\mathbb F$ a field. For $S\subseteq P$ define the function $\nu_S:P\to\mathbb F$ by
\[\nu_S(r)=\begin{cases}0,&r\in S;\\1,&\text{otherwise}.\end{cases}\]
For ease of readability, when $S$ is a singleton $S=\{p\}$, we write $\nu_p$ instead of $\nu_{\{p\}}$.
\end{ntdefinition}
With the notations and definitions above, we are ready to start developing our second rate amplification procedure. We start with the following.
\begin{ntdefinition}[Author Definition39]
Let $P$ be a set, $\mathbb F$ a field, and let $(D^P,L^P)$ be a $(q,\tau,\varepsilon,\rho)$-dual SLR. Let $L^{P^2}$ be as in Definition29. Let $\mathcal R$ be a $(c,s)$-axis evasive partition of $P^2$. We define for every $p\in P^2$ the distribution $(D_{\mathcal R}^{P^2})_p$ as follows. To sample $u$ from $(D_{\mathcal R}^{P^2})_p$:
\begin{enumerate}
\item Sample $g\sim D_p^{P^2}$.
\item If $g=\bot$ return $\bot$; Otherwise, denote $L=\{r\in S_p\mid g(r)\ne0\}$ and proceed as follows.
\item For every $r\in L$ and $w\in(r)$ sample $h_{r,w}\sim D_w^{P^2}$.
\item If there exist $r\in L$ and $w\in(r)$ such that either $h_{r,w}=\bot$ or $h_{r,w}(p)\ne0$ return $\bot$. Otherwise return
\begin{equation}u=g\nu_L+\sum_{r\in L}g(r)\sum_{w\in(r)}\frac{h_{r,w}\nu_w}{h_{r,w}(w)}.\tag{4.3}\end{equation}
\end{enumerate}
Note that, upon reaching Step(4), $u$ is well-defined as $h_{r,w}(w)\ne0$ for all $r\in L$ and $w\in(r)$. We denote the collection of distributions $\{(D_{\mathcal R}^{P^2})_p\mid p\in P^2\}$ by $D_{\mathcal R}^{P^2}$.
\end{ntdefinition}
We start by analyzing the function $u$ that is given by Equation(4.3) above.
\begin{ntclaim}[Author Claim40]
With the notation of Definition39, if $\bot$ is not returned then $u\in F_p$.
\end{ntclaim}
\begin{proof}
As $\bot$ was not returned, for every $r\in L$ and $w\in(r)$ it holds that $h_{r,w}\ne\bot$ and $h_{r,w}(p)=0$. Substituting to Equation(4.3), we get
\[u(p)=g(p)\nu_L(p)=g(p)\ne0,\]
where the second equality holds as $p\notin L$ and the last inequality follows since $g\in\operatorname{supp}(D_p^{P^2})\setminus\{\bot\}$.
\end{proof}
\begin{ntclaim}[Author Claim41]
With the notation of Definition39, if $\bot$ is not returned then $u\in L^{P^2}+L_{\mathcal R}$.
\end{ntclaim}
\begin{proof}
Take $f\in(L^{P^2}+L_{\mathcal R})^\perp$. To prove the claim, it suffices to show that $\langle u,f\rangle=0$. Indeed, this would imply $u\in((L^{P^2}+L_{\mathcal R})^\perp)^\perp=L^{P^2}+L_{\mathcal R}$. As $u\ne\bot$ we have that $g\ne\bot$. Note that
\[\langle g\nu_L,f\rangle=\langle g,f\rangle-\sum_{r\in L}g(r)f(r).\]
Since $g\in\operatorname{supp}(D_p^{P^2})$ we get that $g\in L^{P^2}$. However, $f\in(L^{P^2}+L_{\mathcal R})^\perp\subseteq(L^{P^2})^\perp$, implying $\langle g,f\rangle=0$. Thus,
\begin{equation}\langle g\nu_L,f\rangle=-\sum_{r\in L}g(r)f(r).\tag{4.4}\end{equation}
Now, fix $r\in L$ and $w\in(r)$. By Definition39, as $u\ne\bot$ we have that $h_{r,w}\ne\bot$ and so $h_{r,w}\in L^{P^2}$. However, by the above, $f\in(L^{P^2})^\perp$ and so $\langle h_{r,w},f\rangle=0$. Thus,
\[\langle h_{r,w}\nu_w,f\rangle=\langle h_{r,w},f\rangle-h_{r,w}(w)f(w)=-h_{r,w}(w)f(w).\]
Therefore, for every fixed $r\in L$ one has that
\begin{align}\left\langle\sum_{w\in(r)}\frac{h_{r,w}\nu_w}{h_{r,w}(w)},f\right\rangle&=\sum_{w\in(r)}\left\langle\frac{h_{r,w}\nu_w}{h_{r,w}(w)},f\right\rangle\notag\\&=\sum_{w\in(r)}\frac1{h_{r,w}(w)}\langle h_{r,w}\nu_w,f\rangle\notag\\&=-\sum_{w\in(r)}f(w).\tag{4.5}\end{align}
Now, $f\in(L^{P^2}+L_{\mathcal R})^\perp\subseteq(L_{\mathcal R})^\perp$ whereas $g_{[r]}\in L_{\mathcal R}$, and so
\[0=\langle f,g_{[r]}\rangle=\sum_{w\in[r]}f(w).\]
Substituting this to Equation(4.5), we get
\[\left\langle\sum_{w\in(r)}\frac{h_{r,w}\nu_w}{h_{r,w}(w)},f\right\rangle=f(r).\]
Therefore,
\begin{align*}\left\langle\sum_{r\in L}g(r)\sum_{w\in(r)}\frac{h_{r,w}\nu_w}{h_{r,w}(w)},f\right\rangle&=\sum_{r\in L}g(r)\left\langle\sum_{w\in(r)}\frac{h_{r,w}\nu_w}{h_{r,w}(w)},f\right\rangle\\&=\sum_{r\in L}g(r)f(r).\end{align*}
The above equation together with Equation(4.4) yield $\langle u,f\rangle=0$.
\end{proof}
\begin{ntclaim}[Author Claim42]
With the notation of Definition39, for every $p\in P^2$,
\[\Pr[(D_{\mathcal R}^{P^2})_p=\bot]\leq18csq\tau^2+2cq\varepsilon.\]
\end{ntclaim}
\begin{proof}
First, the probability that $g=\bot$ is bounded by $\varepsilon$. Similarly, the probability that for any specific $r\in L$ and $w\in(r)$, $h_{r,w}=\bot$ is bounded by $\varepsilon$. Thus, by the union bound, and since $|L|\leq2q-1$ and $|(r)|\leq c$, we have that expect with probability $(1+(2q-1)c)\varepsilon\leq2qc\varepsilon$, the sampling process above will result in $u\ne\bot$.
To complete the analysis, we turn to bound the probability that $h_{r,w}(p)=0$ for some $r\in L$ and $w\in(r)$. Let $L=\{r_1,\ldots,r_{|L|}\}$. While the random variables in $L$ may be dependent, marginally, it holds that for every $i\in[|L|]$ and every fixed $r\in S_p$, $\Pr[r_i=r]\leq\tau$. With this notation, by Definition39, $(D_{\mathcal R}^{P^2})_p=\bot$ only if there exist $i\in[|L|]$ and $w\in(r_i)$ such that $h_{r_i,w}(p)\ne0$.
For a fixed $r\in S_p$ define the event $E_r$ in which there exists $w\in(r)$ such that $h_{r,w}(p)\ne0$, (when conditioned on $h_{r,w}\ne\bot$). Note that this event is with respect to the randomness of sampling $h_r=\{h_{r,w}\mid w\in(r)\}$ whereas $r$ is fixed. By the union bound,
\[\Pr_{h_r}[E_r]\leq\sum_{w\in(r)}\Pr_{h_{r,w}}[h_{r,w}(p)\ne0\mid h_{r,w}\ne\bot].\]
Observe first that $w\ne p$. Indeed, as $r\in S_p$, both $r$ and $p$ are on some common axis-parallel line $\ell\in\mathcal X$. Thus, $w=p$ would imply $|[r]\cap\ell|\geq2$ which stands in contradiction to the definition of axis-evasiveness. Consider $w\in(r)\setminus S_p$. As $w\ne p$ we have that $\operatorname{dist}(w,p)=2$. By Lemma32, as $h_{r,w}\sim D_w^{P^2}$ we have that
\[\Pr_{h_{r,w}}[h_{r,w}(p)\ne0\mid h_{r,w}\ne\bot]\leq\tau^2.\]
If, on the other hand, $w\in(r)\cap S_p$ then $\operatorname{dist}(w,p)=1$, and Lemma32 then implies that
\[\Pr_{h_{r,w}}[h_{r,w}(p)\ne0\mid h_{r,w}\ne\bot]\leq\tau.\]
As $|(r)|\leq c$ we conclude that
\[\Pr_{h_r}[E_r]\leq c\tau^2+\tau|(r)\cap S_p|.\]
Fix $i\in[|L|]$ and consider the random variable $r_i$. The above equation, together with $|(r_i)|\leq c$, yields
\begin{align}\Pr_{r_i,h_{r_i}}[E_{r_i}]&\leq\Pr_{r_i,h_{r_i}}[E_{r_i}\mid(r_i)\cap S_p=\varnothing]\notag\\&\quad+\Pr_{r_i,h_{r_i}}[E_{r_i}\mid(r_i)\cap S_p\ne\varnothing]\Pr_{r_i}[(r_i)\cap S_p\ne\varnothing]\notag\\&\leq c\tau^2+(c\tau^2+c\tau)\Pr_{r_i}[(r_i)\cap S_p\ne\varnothing].\tag{4.6}\end{align}
Consider now the set $B=\{r\in S_p\mid(r)\cap S_p\ne\varnothing\}$. As $\mathcal R$ is $(c,s)$-axis evasive, Claim36 implies that $|B|\leq4s$, and so
\[\Pr_{r_i}[(r_i)\cap S_p\ne\varnothing]=\Pr[r_i\in B]\leq4s\tau.\]
Substituting to Equation(4.6), we get $\Pr[E_i]\leq9cs\tau^2$. The proof then follows by taking the union bound over all $i\in[|L|]$ as, indeed, $|L|=2q-1$.
\end{proof}
\begin{ntclaim}[Author Claim43]
With the notation of Definition39, for every pair of distinct $p,r\in P^2$,
\[\Pr_{u\sim(D_{\mathcal R}^{P^2})_p}[u(r)\ne0\mid u\ne\bot]\leq10csq\tau^2.\]
\end{ntclaim}
\begin{proof}
By Equation(4.3), $u$ is a linear combination of the (sampled) functions $g\nu_L,\{h_{r,w}\nu_w\}$. To prove the claim, we will show that, with high probability, all these functions evaluate to $0$ at the point $r$, implying $u(r)=0$. We start by bounding $\Pr[(g\nu_L)(r)\ne0]$. To this end, consider two cases. First, if $r\in P^2\setminus S_p$ then, as $L\subseteq S_p$, we have that $\nu_L(r)=1$ and so in such case
\begin{equation}\Pr[(g\nu_L)(r)\ne0]=\Pr[g(r)\ne0]\leq\tau^2,\tag{4.7}\end{equation}
where the last inequality follows by Lemma32 and since $\operatorname{dist}(r,p)=2$ per our assumption $r\notin S_p$ and since $r\ne p$. If, on the other hand, $r\in S_p$ then, by the definition of $L$,
\[g(r)\ne0\Longrightarrow r\in L\Longrightarrow\nu_L(r)=0,\]
and so in this case $(g\nu_L)(r)=0$.
Let $L=\{r_1,\ldots,r_{|L|}\}$. Consider a fixed $i\in[|L|]$ and denote $(r_i)=\{w_{i,1},\ldots,w_{i,b}\}$, where $b\leq c$. Fix $j\in[b]$. We turn to bound $\Pr[(h_{r_i,w_{i,j}}\nu_{w_{i,j}})(r)\ne0]$. First note that
\begin{equation}\Pr[(h_{r_i,w_{i,j}}\nu_{w_{i,j}})(r)\ne0\mid(r_i)\cap S_r=\varnothing]\leq\tau^2.\tag{4.8}\end{equation}
Indeed, conditioned on the event $(r_i)\cap S_r=\varnothing$, either $w_{i,j}=r$ or $\operatorname{dist}(w_{i,j},r)=2$. In the first case,
\[(h_{r_i,w_{i,j}}\nu_{w_{i,j}})(r)=h_{r_i,r}(r)\nu_r(r)=0.\]
In the second case, the bound follows by Lemma32. Second, note that
\begin{equation}\Pr[(h_{r_i,w_{i,j}}\nu_{w_{i,j}})(r)\ne0\mid(r_i)\cap S_r\ne\varnothing]\leq\tau.\tag{4.9}\end{equation}
Indeed, as before, we may only consider the case $r\ne w_{i,j}$ and then observe that $\operatorname{dist}(r,w_{i,j})=1$ and invoke Lemma32. Now, let $B=\{v\in S_p\mid(v)\cap S_r\ne\varnothing\}$. By Claim36, and since $\mathcal R$ is $(c,s)$-axis evasive, $|B|\leq4s$. Recall that $\Pr[r_i=v]\leq\tau$ for every fixed $v\in S_p$, and so
\begin{equation}\Pr[(r_i)\cap S_r\ne\varnothing]=\Pr[r_i\in B]\leq4s\tau.\tag{4.10}\end{equation}
By Equations(4.8),(4.9),(4.10) we get
\[\Pr[(h_{r_i,w_{i,j}}\nu_{w_{i,j}})(r)\ne0]\leq\tau^2+4s\tau^2\leq5s\tau^2.\]
The proof then follows by the union bound over all $i\in[|L|]$ and $j\in[|(w_i)|]$.
\end{proof}
\begin{ntdefinition}[Author Definition44]
Let $P$ be a set, $\mathbb F$ a field, and let $(D^P,L^P)$ be a $(q,\tau,\varepsilon,\rho)$-dual SLR. Let $L^{P^2}$ be as in Definition29. Let $\mathcal R$ be a $(c,s)$-axis evasive partition of $P^2$ and let $D_{\mathcal R}^{P^2}$ be as in Definition39. We denote by $(D^P,L^P)^2_{\mathcal R}$ the pair $(D_{\mathcal R}^{P^2},L^{P^2}+L_{\mathcal R})$.
\end{ntdefinition}
\begin{ntproposition}[Author Proposition45]
Let $P$ be a set, $\mathbb F$ a field, and let $(D^P,L^P)$ be a $(q,\tau,\varepsilon,\rho)$-dual SLR. Let $\mathcal R$ be a $(c,s)$-axis evasive partition of $P^2$ that consists of $t$ parts. Then, $(D^P,L^P)^2_{\mathcal R}$ is a $(q_{\mathcal R},\tau_{\mathcal R},\varepsilon_{\mathcal R},\rho_{\mathcal R})$-dual SLR with
\begin{align*}q_{\mathcal R}&\leq2cq^3\\\tau_{\mathcal R}&\leq10csq\tau^2\\\varepsilon_{\mathcal R}&\leq18csq\tau^2+2cq\varepsilon\\\rho_{\mathcal R}&\geq1-(1-\rho)^2-\frac t{|P|^2}.\end{align*}
\end{ntproposition}
\begin{proof}
Claim40 implies that for every $p\in P^2$, $\operatorname{supp}((D_{\mathcal R}^{P^2})_p)\subseteq F_p\cup\{\bot\}$. To bound $q_{\mathcal R}$, note that by Equation(4.3),
\[|u|\leq|g\nu_L|+\sum_{r\in L}\sum_{w\in(r)}|h_{r,w}\nu_w|.\]
Now, $|g\nu_L|\leq|g|\leq q^2$ and $|h_{r,w}\nu_w|\leq|h_{r,w}|\leq q^2$. Hence, $|u|\leq q^2+|L|cq^2\leq2cq^3$. The stated bounds on $\tau_{\mathcal R}$ and $\varepsilon_{\mathcal R}$ readily follows by Claim43 and Claim42, respectively. As for the rate, we have that
\begin{align*}\dim(L^{P^2}+L_{\mathcal R})&\leq\dim(L^{P^2})+\dim(L_{\mathcal R})\\&\leq(1-\rho)^2|P|^2+t,\end{align*}
where the second inequality follows by Proposition34 and since $\mathcal R$ consists of $t$ parts, implying $|L_{\mathcal R}|=t$.
\end{proof}
